% from ega4-2-fr.tex
\oldpage[IV-2]{119}

canoniques, l'homomorphisme canonique
$\rho_{X'/Z'}:\mathcal F'\to i'_*(\mathcal F'|V')$ est tel que
$g_*(\rho_{X'/Z'})$ soit l'homomorphisme canonique
\[
\rho_{X/Z}:\mathcal H^0_{X/Z}(\mathcal F)
\longrightarrow i_*(\mathcal H^0_{X/Z}(\mathcal F)|V)
\]
(II, 1.4.2), compte tenu de (5.10.11.1). Comme
$\mathcal H^0_{X/Z}(\mathcal F)$ est $Z$-clos (5.9.11), $\rho_{X/Z}$ est un
isomorphisme, donc il en est de même de $\rho_{X'/Z'}$. Comme $X'-Z'$ est
intersection de la famille filtrante des $V'_\alpha=g^{-1}(V_\alpha)$, où
$V_\alpha$ parcourt la famille filtrante des ouverts contenant $X-Z$, on en
déduit que $\mathcal F'$ est $Z'$-clos lorsque $X'$ est localement noethérien,
en vertu de (5.9.1).

\begin{env}[5.10.13]
\label{IV.2.s5.10.13-fr}
Nous allons maintenant appliquer les résultats qui précèdent au cas où $Z$ est
un des ensembles $Z^{(n)}(X)$ (ou simplement $Z^{(n)}$), défini comme
l'ensemble des $x\in X$ tels que $\dim(\mathcal O_x)\geq n$; il est clair que
$Z^{(n)}$ est stable par spécialisation; pour qu'une partie fermée $T$ de $X$
soit contenue dans $Z^{(n)}$, il faut et il suffit que
$\operatorname{codim}(T,X)\geq n$. Nous nous intéresserons ici au cas $n=2$.
\end{env}

\begin{proposition}[5.10.14]
\label{IV.2.s5.10.14-fr}
Soient $X$ un préschéma localement noethérien, $\mathcal F$ un
$\mathcal O_X$-Module cohérent de support égal à $X$.
\begin{enumerate}
\item[(i)] Pour que $\mathcal F$ vérifie $(S_1)$, il faut et il suffit qu'il
soit $Z^{(1)}$-pur.
\item[(ii)] Pour que $\mathcal F$ vérifie $(S_2)$, il faut et il suffit qu'il
soit $Z^{(2)}$-clos et $Z^{(1)}$-pur, ou encore qu'il soit $Z^{(2)}$-clos et
n'ait pas de cycle premier associé de codimension 1.
\end{enumerate}
\end{proposition}

(i) Dire que $\mathcal F$ possède la propriété $(S_1)$ signifie que
$\mathcal F$ n'a pas de cycle premier associé immergé (5.7.5), ou encore que
pour tout $x\in\operatorname{Ass}(\mathcal F)$, on a
$\dim(\mathcal F_x)=0$ (3.1.4), autrement dit (5.1.12.1)
$\dim(\mathcal O_x)=0$; mais cela équivaut à dire que
$\operatorname{Ass}(\mathcal F)$ ne rencontre pas $Z^{(1)}$, et la conclusion
résulte de (5.10.2).

(ii) Dire que $\mathcal F$ est $Z^{(2)}$-clos signifie que
$\operatorname{prof}_{Z^{(2)}}(\mathcal F)\geq2$, ou encore que, pour tout
$x\in X$, la relation $\dim(\mathcal O_x)\geq2$ entraîne
$\operatorname{prof}(\mathcal F_x)\geq2$; cela montre que la propriété $(S_2)$
entraîne que $\mathcal F$ est $Z^{(2)}$-clos; elle entraîne en outre que
$\mathcal F$ vérifie $(S_1)$, donc n'a pas de cycle premier associé immergé
(5.7.5), et comme $\operatorname{Supp}(\mathcal F)=X$, cela signifie encore
que tous les cycles premiers associés à $\mathcal F$ sont de codimension 0.
Inversement, supposons que $\mathcal F$ soit $Z^{(2)}$-clos et n'ait pas de
cycle premier associé de codimension 1; pour voir que $\mathcal F$ vérifie
$(S_2)$, il reste à montrer que si $x\in X$ est tel que
$\dim(\mathcal F_x)=1$ (ou, ce qui revient au même,
$\dim(\mathcal O_x)=1$), alors on a $\operatorname{prof}(\mathcal F_x)=1$;
mais par hypothèse la relation $\dim(\mathcal O_x)=1$ entraîne
$x\notin\operatorname{Ass}(\mathcal F)$, et cette dernière relation équivaut à
$\operatorname{prof}(\mathcal F_x)\neq0$, c'est-à-dire ici à
$\operatorname{prof}(\mathcal F_x)=1$. Si $\mathcal F$ est $Z^{(1)}$-pur,
donc vérifie $(S_1)$, on a remarqué plus haut qu'en vertu de la relation
$\operatorname{Supp}(\mathcal F)=X$, tous les cycles premiers associés à
$\mathcal F$ sont de codimension 0, donc ce qui précède s'applique.

On notera qu'il peut se faire que $\mathcal F$ soit $Z^{(2)}$-clos et ne
vérifie pas $(S_1)$: c'est le cas par exemple lorsque $X$ est de dimension 1
(car alors $Z^{(2)}=\varnothing$, et tout $\mathcal O_X$-Module est
$Z^{(2)}$-clos) et a des cycles premiers associés immergés.

Rappelons qu'au chapitre III, dans l'étude de la cohomologie locale, on donne
une caractérisation cohomologique de la propriété $(S_n)$ pour tout $n\geq1$,
généralisant (5.10.14).

\begin{corollary}[5.10.15]
\label{IV.2.s5.10.15-fr}
Soient $X$ un préschéma localement noethérien, $\mathcal F$ un
$\mathcal O_X$-Module cohérent de support $X$. On suppose que $\mathcal F$
n'a pas de cycles premiers associés de codimension 1 et que
$\mathcal F'=\mathcal H^0_{X/Z^{(2)}}(\mathcal F)$ soit cohérent. Alors:
\begin{enumerate}
\item[(i)] $\mathcal F'$ vérifie la propriété $(S_2)$.

